\begin{abstract}
We examined whether imagined left-hand, right-hand, feet and tongue movements have different topological signatures in EEG sensor-dependence networks. Training recordings from all nine participants in BCI Competition IV data set 2a were analysed. After original artifact exclusions, 2,328 of 2,592 trials remained; balancing the four conditions within each participant selected 2,224 trials. Trial-pooled Welch coherence in mu (8--12 Hz bins) and beta (13--30 Hz bins) bands was transformed to channel dissimilarity, followed by weighted flag filtrations, $H_0$/$H_1$ persistence and standard stable ranks. Individual movement differences were visible, particularly in A01, but their ordering and size varied across people. Exploratory permutation-calibrated Friedman tests of finite total persistence did not detect a group-level movement effect after Holm correction (adjusted $p\geq0.9816$). A common sample size of 49 trials per condition gave the same inferential conclusion. The findings demonstrate a reproducible approach to frequency-specific EEG topology, while providing insufficient evidence for a consistent movement effect in the tested scalar summaries. They do not establish equality of the full networks or barcodes.
\end{abstract}
